\section*{Capítulo \ref{ch:g9:earth-resources} --- Los recursos de la Tierra}

\begin{solution}{exo:g9:earth-resources:1}
(a) \omterm{def:g9:earth-resources:aggregate}{Áridos}. (b) \omterm{def:g9:earth-resources:reserve}{Reservas}. (c) \omterm{def:g9:earth-resources:ore-deposit}{Yacimiento mineral}. (d) \omterm{def:g9:earth-resources:ore-deposit}{Factor de concentración}.
\end{solution}

\begin{solution}{exo:g9:earth-resources:2}
(a) En las \omterm{def:g3:stones-pebbles-sand:pebble}{gravas} de los \omterm{def:g1:landscapes-around-us:water}{ríos}. (b) En las grietas, a partir de agua
caliente. (c) En el \omterm{def:g7:volcanoes:magma}{magma}. (d) Es el primer paso de un yacimiento de
\omterm{def:g3:stones-pebbles-sand:pebble}{gravas} fluviales: el oro liberado lo transporta después el \omterm{def:g1:landscapes-around-us:water}{río}.
\end{solution}

\begin{solution}{exo:g9:earth-resources:3}
(b), (d), (a), (e), (c).
\end{solution}

\begin{solution}{exo:g9:earth-resources:4}
(a) No renovable: tarda millones de años en formarse. (b) Renovable: la
\omterm{def:g4:clouds-rain-snow:precipitation}{lluvia} repone lo que se extrae. (c) No renovable: un \omterm{def:g9:earth-resources:ore-deposit}{yacimiento mineral} se
forma a lo largo de tiempos muy largos. (d) No renovable a ese ritmo: el
\omterm{def:g6:water-underground:water-table}{nivel freático} baja y los \omterm{def:g6:water-underground:spring}{pozos} acaban secos.
\end{solution}

\begin{solution}{exo:g9:earth-resources:5}
(a) $4 \div 0.007 \approx 571$, del orden de $10^3$ (entre $10^2$ y
$10^3$, más cerca de $10^3$). (b) $0.5 \div 0.0002 = 2500$, del orden de
$10^3$. (c) $20
\div 5.6 \approx 3.6$, unas 4 veces: entre $10^0$ y $10^1$.
(d) El hierro ya es muy abundante en la \omterm{def:g6:inside-the-earth:layers}{corteza}: basta una mena solo unas
pocas veces más rica que la \omterm{def:g3:stones-pebbles-sand:rock}{roca} corriente.
\end{solution}

\begin{solution}{exo:g9:earth-resources:6}
(a) $19.3 \div 2.5 \approx 7.7$ veces. (b) Oro: $19.3 \times 0.01 =
\qty{0.193}{g}$; \omterm{def:g3:stones-pebbles-sand:rock}{roca}: $2.5 \times 0.01 = \qty{0.025}{g}$. (c) Cuando todo
el lecho se mueve, los granos de oro, pesados, se hunden entre la \omterm{def:g3:stones-pebbles-sand:sand}{arena} y
los \omterm{def:g3:stones-pebbles-sand:pebble}{cantos}, más ligeros, hasta la \omterm{def:g3:stones-pebbles-sand:rock}{roca} madre sólida, donde se detienen.
\end{solution}

\begin{solution}{exo:g9:earth-resources:7}
(a) El oro, el estaño y el cromo. (b) Entre $10^1$ y $10^2$, cerca de
$10^2$. (c) $3000 \div 3.7 \approx 800$ veces. (d) El cromo (o el estaño, o
el oro): cuanto mayor es el \omterm{def:g9:earth-resources:ore-deposit}{factor de concentración}, más escasos son los
lugares donde la naturaleza ha concentrado el metal lo suficiente.
\end{solution}

\begin{solution}{exo:g9:earth-resources:8}
(a) $18 \times 365 = \qty{6570}{kg}$, unas \qty{6.6}{t}. (b) $5 \times
10^{10}$ toneladas. (c) $25000 \div 18 \approx 1389$ días, casi cuatro
años. (d) La \omterm{def:g3:stones-pebbles-sand:sand}{arena} sacada del \omterm{def:g1:landscapes-around-us:water}{río} ya no viaja aguas abajo hasta la \omterm{def:g1:landscapes-around-us:water}{costa};
las olas siguen llevándose \omterm{def:g3:stones-pebbles-sand:sand}{arena} de la playa, pero llega menos para
reponerla.
\end{solution}

\begin{solution}{exo:g9:earth-resources:9}
(a) De diminutas plantas y animales del \omterm{def:g2:water-on-earth:ocean}{mar}, como las diatomeas,
enterrados en el \omterm{def:g3:stones-pebbles-sand:sand}{barro} del fondo marino. (b) Entre unos 60 y
\qty{120}{\degreeCelsius}. (c) Son más ligeros que el agua, que llena los
\omterm{def:g6:water-underground:porosity}{poros} de las \omterm{def:g3:stones-pebbles-sand:rock}{rocas}. (d) Una capa \omterm{def:g6:water-underground:permeability}{impermeable}, como la lutita; se acumulan
bajo ella, a menudo en lo alto de un pliegue abombado, el gas encima del
petróleo.
\end{solution}

\begin{solution}{exo:g9:earth-resources:10}
(a) $2.9 \times 10^7$ toneladas. (b) $2.9 \times 10^7 \div 100 = 2.9 \times
10^5$ toneladas al año, aproximadamente. (c) La mena pesa 100 veces más que
el cobre que contiene: $2.9 \times 10^7 \times 100 = 2.9 \times 10^9$
toneladas. (d) Cada tonelada de mena solo da \qty{10}{kg} de cobre, así
que hay que excavar miles de millones de toneladas de \omterm{def:g3:stones-pebbles-sand:rock}{roca}.
\end{solution}

\begin{solution}{exo:g9:earth-resources:11}
(a) A 23~millones de toneladas al año habrían durado unos $280 \div 23
\approx 12$ años, hasta principios de la década de 1980; sin embargo, se
extrajeron 697~millones de toneladas, mucho más que esas \omterm{def:g9:earth-resources:reserve}{reservas}.
(b) $980 \div 280 = 3.5$ veces. (c) $980 \div
23 \approx 43$ años. (d) Por ejemplo: se descubren nuevos yacimientos;
unas técnicas mejores o unos precios más altos hacen rentables
yacimientos más pobres.
\end{solution}

\begin{solution}{exo:g9:earth-resources:12}
(a) $5 \times 10^1$ años. (b) $2 \times 10^3 \div 50 = 40$ veces. (c) $3
\times 10^8 \div (5 \times 10^1) = 6 \times 10^6$ veces. (d) El carbón
tardó millones de veces más en formarse que lo que tardamos en usarlo: la
naturaleza no lo renueva a nuestra escala de tiempo.
\end{solution}

\begin{solution}{pb:g9:earth-resources:1}
\textbf{1.} $100 \times 1000 = \num{100000}$~kg $= 1 \times 10^5$~kg.

\textbf{2.} $2000 \times 300 = \num{600000}$~t $= 6 \times 10^5$~t.

\textbf{3.} $6 \times 10^5 \div 10 = 6 \times 10^4$~t al año.

\textbf{4.} $0.40 \times 6 \times 10^4 = 2.4 \times 10^4$~t de \omterm{def:g3:stones-pebbles-sand:sand}{arena} al
año.

\textbf{5.} $1200 \times 500 \times 5 = 3 \times 10^6$~m$^3$.

\textbf{6.} Dos tercios: $2 \times 10^6$~m$^3$.

\textbf{7.} $2 \times 10^6 \times 1.6 = 3.2 \times 10^6$~t.

\textbf{8.} $\num{24000} + \num{136000} = \num{160000}$~t $= 1.6 \times
10^5$~t al año.

\textbf{9.} $\num{160000} \div 2000 = 80$ veces más.

\textbf{10.} No: la gravera se vacía 80 veces más deprisa de lo que el \omterm{def:g1:landscapes-around-us:water}{río}
renueva la \omterm{def:g3:stones-pebbles-sand:sand}{arena}. A la escala de la ciudad, esta \omterm{def:g3:stones-pebbles-sand:sand}{arena} es un recurso
limitado.

\textbf{11.} $\num{24000} \div 4 = \num{6000}$~t de \omterm{def:g3:stones-pebbles-sand:sand}{arena} al año.

\textbf{12.} $3.2 \times 10^6 \div (1.6 \times 10^5) = 20$: \textbf{20 años}.
\end{solution}
